% comparison.tex - Comparison layout macros for Paperforge
% Requires: _colors.tex, _styles.tex loaded first
%
% Macros provided:
%   \sidebyside{left title}{left content}{right title}{right content}
%   \beforeafter{before title}{before content}{after title}{after content}
%   \proscons{pros list}{cons list}
%   \comparisonthree{col1}{col2}{col3}
%
% Design: Green tint for positive/before, Orange tint for negative/after

% ============================================================================
% SIDE-BY-SIDE COMPARISON
% ============================================================================
% Usage: \sidebyside{Left Title}{Left content}{Right Title}{Right content}
\newcommand{\sidebyside}[4]{%
\begin{center}
\begin{tikzpicture}
    % Left column
    \node[diagram compare left, minimum height=3cm] (left) at (0,0) {%
        \textbf{#1}\\[4pt]
        \begin{minipage}{3.5cm}
        #2
        \end{minipage}
    };
    % Right column
    \node[diagram compare right, minimum height=3cm] (right) at (5.5,0) {%
        \textbf{#3}\\[4pt]
        \begin{minipage}{3.5cm}
        #4
        \end{minipage}
    };
    % Center divider
    \draw[diagram divider] (2.75, 1.7) -- (2.75, -1.7);
\end{tikzpicture}
\end{center}
}

% ============================================================================
% BEFORE/AFTER COMPARISON
% ============================================================================
% Usage: \beforeafter{Before Title}{Before content}{After Title}{After content}
\newcommand{\beforeafter}[4]{%
\begin{center}
\begin{tikzpicture}
    % Before box
    \node[diagram compare left, minimum height=2.5cm, minimum width=4.5cm] (before) at (0,0) {%
        \textbf{#1}\\[4pt]
        \begin{minipage}{4cm}
        #2
        \end{minipage}
    };
    % Arrow
    \node at (3.5, 0) {\Large\color{diagramPrimary}$\Rightarrow$};
    % After box
    \node[diagram compare right, minimum height=2.5cm, minimum width=4.5cm] (after) at (7,0) {%
        \textbf{#3}\\[4pt]
        \begin{minipage}{4cm}
        #4
        \end{minipage}
    };
\end{tikzpicture}
\end{center}
}

% ============================================================================
% PROS AND CONS
% ============================================================================
% Usage: \proscons{Pro item 1\\Pro item 2}{Con item 1\\Con item 2}
\newcommand{\proscons}[2]{%
\begin{center}
\begin{tikzpicture}
    % Pros column (green)
    \node[rectangle, rounded corners=4pt, draw=diagramSuccess!50, fill=diagramPositive,
          minimum width=5cm, minimum height=3cm, align=left, inner sep=10pt] (pros) at (0,0) {%
        \textcolor{diagramSuccess}{\textbf{\large $\checkmark$ Pros}}\\[6pt]
        \begin{minipage}{4.5cm}
        \small #1
        \end{minipage}
    };
    % Cons column (orange/red)
    \node[rectangle, rounded corners=4pt, draw=diagramError!50, fill=diagramError!10,
          minimum width=5cm, minimum height=3cm, align=left, inner sep=10pt] (cons) at (6,0) {%
        \textcolor{diagramError}{\textbf{\large $\times$ Cons}}\\[6pt]
        \begin{minipage}{4.5cm}
        \small #2
        \end{minipage}
    };
\end{tikzpicture}
\end{center}
}

% ============================================================================
% THREE-COLUMN COMPARISON
% ============================================================================
% Usage: \comparisonthree{Title1}{Content1}{Title2}{Content2}{Title3}{Content3}
\newcommand{\comparisonthree}[6]{%
\begin{center}
\begin{tikzpicture}
    % Column 1
    \node[diagram card, minimum width=3.5cm, minimum height=3cm] (c1) at (0,0) {%
        \textbf{#1}\\[4pt]
        \begin{minipage}{3cm}
        \small #2
        \end{minipage}
    };
    % Column 2
    \node[diagram card, minimum width=3.5cm, minimum height=3cm] (c2) at (4.5,0) {%
        \textbf{#3}\\[4pt]
        \begin{minipage}{3cm}
        \small #4
        \end{minipage}
    };
    % Column 3
    \node[diagram card, minimum width=3.5cm, minimum height=3cm] (c3) at (9,0) {%
        \textbf{#5}\\[4pt]
        \begin{minipage}{3cm}
        \small #6
        \end{minipage}
    };
\end{tikzpicture}
\end{center}
}

% ============================================================================
% FEATURE COMPARISON TABLE (Visual)
% ============================================================================
% Usage: \featurecompare{Feature}{Option A}{Option B}
%        Call multiple times, wrap in tikzpicture
\newcommand{\featurerow}[3]{%
    \node[anchor=west, font=\small] at (0, \featureY) {#1};
    \node[anchor=center] at (5, \featureY) {#2};
    \node[anchor=center] at (8, \featureY) {#3};
    \pgfmathsetmacro{\featureY}{\featureY - 0.6}
}

% ============================================================================
% DARK THEME SIDE-BY-SIDE
% ============================================================================
\newcommand{\sidebysidedark}[4]{%
\begin{center}
\begin{tikzpicture}
    % Left column
    \node[diagram card dark, minimum height=3cm, minimum width=4.5cm] (left) at (0,0) {%
        \textbf{#1}\\[4pt]
        \begin{minipage}{4cm}
        #2
        \end{minipage}
    };
    % Right column
    \node[diagram card dark, minimum height=3cm, minimum width=4.5cm] (right) at (5.5,0) {%
        \textbf{#3}\\[4pt]
        \begin{minipage}{4cm}
        #4
        \end{minipage}
    };
    % Center divider
    \draw[draw=diagramDarkBorder, line width=1pt] (2.75, 1.7) -- (2.75, -1.7);
\end{tikzpicture}
\end{center}
}

% ============================================================================
% VERSUS COMPARISON
% ============================================================================
% Usage: \versus{Option A}{Option B}
\newcommand{\versus}[2]{%
\begin{center}
\begin{tikzpicture}
    % Option A
    \node[diagram box primary, minimum width=4cm, minimum height=1.5cm] (a) at (0,0) {%
        \textbf{#1}
    };
    % VS badge
    \node[circle, fill=diagramAccent, minimum size=1cm, font=\small\bfseries] at (2.75, 0) {VS};
    % Option B
    \node[diagram box secondary, minimum width=4cm, minimum height=1.5cm] (b) at (5.5,0) {%
        \textbf{#2}
    };
\end{tikzpicture}
\end{center}
}
